\documentclass[12pt]{ximera}
\title{Exam 3 (2023)}
\begin{document}
\begin{abstract}
An archived exam on Chapter 3 and the start of Chapter 4: lone-divider, lone-chooser on a three-topping pizza, sealed bids with assets and chores, markers with point-valued candy, and a Hamilton apportionment.
\end{abstract}
\maketitle

\section*{Exam 3 (2023)}

\begin{problem}
Four sisters divide their parents' estate by the lone-divider method. Their values for
the four shares:
\begin{center}
\begin{tabular}{r|rrrr}
 & Allie & Britney & Cass & Dorothy\\\hline
$s_1$ & 200 & 240 & 230 & 250\\
$s_2$ & 280 & 300 & 270 & 250\\
$s_3$ & 230 & 200 & 260 & 250\\
$s_4$ & 290 & 260 & 240 & 250
\end{tabular}
\end{center}

\begin{prompt}
\textbf{(a)} Who was the divider?
\begin{multipleChoice}
\choice{Allie}
\choice{Britney}
\choice{Cass}
\choice[correct]{Dorothy}
\end{multipleChoice}
\begin{feedback}
Only Dorothy values all shares equally.
\end{feedback}
\end{prompt}

\begin{prompt}
\textbf{(b)} Which shares does Cass find acceptable?
\begin{selectAll}
\choice{$s_1$}
\choice[correct]{$s_2$}
\choice[correct]{$s_3$}
\choice{$s_4$}
\end{selectAll}
\begin{feedback}
Each sister values the estate at $1000$, so a fair share is at least $250$. Allie and
Britney both accept $\set{s_2,s_4}$; Cass accepts $\set{s_2,s_3}$.
\end{feedback}
\end{prompt}

\begin{prompt}
\textbf{(c)} In every fair division, which shares do Cass and Dorothy get?

Cass: share number $\answer{3}$ \quad Dorothy: share number $\answer{1}$
\begin{feedback}
Allie and Britney need $s_2$ and $s_4$ between them, which forces Cass to take $s_3$ and
leaves $s_1$ for Dorothy. The two fair divisions differ only in whether Allie takes
$s_2$ or $s_4$.
\end{feedback}
\end{prompt}
\end{problem}

\begin{problem}
Three friends share a pizza with equal parts cheese, meatball, and BBQ chicken, using
the lone-chooser method.
\begin{itemize}
\item Clark likes meatball twice as much as cheese, and BBQ chicken three times as much
      as cheese.
\item Lois likes meatball and BBQ chicken equally, and likes cheese as much as the
      other two together.
\item Lex likes meatball and BBQ chicken equally and gives cheese no value.
\end{itemize}
The steps: (1) Clark cuts two shares of equal value to him, choosing cheese $+$ meatball
and BBQ chicken. (2) Lois picks the share she prefers, and Clark keeps the other. (3)
Each cuts their share into three sub-shares. (4) Lex picks one sub-share from each.

\begin{prompt}
\textbf{(a)} Clark values cheese, meatball, and BBQ at $\frac16$, $\frac26$, $\frac36$.
What values do Lois and Lex give them?

Lois --- cheese $\answer{1/2}$, meatball $\answer{1/4}$, BBQ $\answer{1/4}$

\smallskip
Lex --- cheese $\answer{0}$, meatball $\answer{1/2}$, BBQ $\answer{1/2}$
\begin{feedback}
If Lois values meatball and BBQ at $y$ each, cheese is $2y$, so $4y=1$.
\end{feedback}
\end{prompt}

\begin{prompt}
\textbf{(b)} Which share does Lois pick in step 2?
\begin{multipleChoice}
\choice[correct]{Cheese $+$ meatball}
\choice{BBQ chicken}
\end{multipleChoice}
\begin{feedback}
To Lois, cheese $+$ meatball is worth $\frac12+\frac14=\frac34$.
\end{feedback}
\end{prompt}

\begin{prompt}
\textbf{(c)} In the simplest subdivision, each sub-share has one topping: Lois cuts
$\frac C2$, $\frac C2$, $M$, and Clark cuts BBQ into thirds. Lex takes the meatball
and one third of the BBQ. What are the final values, each in its owner's view?

Lex $\answer{2/3}$ \quad Lois $\answer{1/2}$ \quad Clark $\answer{1/3}$
\begin{feedback}
Lex: $\frac12+\frac13\cdot\frac12=\frac23$. Lois keeps both halves of the cheese, worth
$\frac12$ to her. Clark keeps two thirds of the BBQ, worth $\frac23\cdot\frac36=\frac13$.
Everyone gets at least $\frac13$, and Lex does especially well.
\end{feedback}
\end{prompt}
\end{problem}

\begin{problem}
Brothers Henry, Tom, and Jeff divide a real estate empire by sealed bids (in millions of
dollars).
\begin{center}
\begin{tabular}{r|rrr}
 & H & T & J\\\hline
Apartments & 24 & 21 & 25\\
Hotels & 36 & 39 & 32\\
Office towers & 21 & 24 & 23\\
Mall projects & 24 & 21 & 25
\end{tabular}
\end{center}

\begin{prompt}
\textbf{(a)} Who gets each asset?

Apartments $\answer{J}$ \quad Hotels $\answer{T}$ \quad Offices $\answer{T}$ \quad
Malls $\answer{J}$
\end{prompt}

\begin{prompt}
\textbf{(b)} Initial settlement (positive = pays): H $\answer{-35}$ \quad T $\answer{28}$
\quad J $\answer{15}$
\begin{feedback}
Each brother's bids total $105$, so every fair share is $35$. Tom's assets are worth
$63$ and Jeff's $50$ to them; Henry has nothing.
\end{feedback}
\end{prompt}

\begin{prompt}
\textbf{(c)} Surplus: $\answer{8}$ million
\begin{feedback}
$28+15-35=8$, or about $\$2.67$ million per brother.
\end{feedback}
\end{prompt}
\end{problem}

\begin{problem}
Housemates A, B, C use sealed bids to split a week of chores; the amounts are what each
would charge to do the chore.
\begin{center}
\begin{tabular}{r|rrr}
 & A & B & C\\\hline
Trash & 39 & 42 & 36\\
Vacuum & 25 & 24 & 29\\
Dishes & 26 & 24 & 25
\end{tabular}
\end{center}

\begin{prompt}
\textbf{(a)} Who does each chore?

Trash $\answer{C}$ \quad Vacuum $\answer{B}$ \quad Dishes $\answer{B}$
\begin{feedback}
Each chore goes to the lowest bidder; A does nothing.
\end{feedback}
\end{prompt}

\begin{prompt}
\textbf{(b)} Initial settlement (positive = pays): A $\answer{30}$ \quad B $\answer{-18}$
\quad C $\answer{-6}$
\begin{feedback}
Everyone's total is $90$, so a fair share of the work is $30$. A does $0$ and pays $30$;
B does $48$ and receives $18$; C does $36$ and receives $6$.
\end{feedback}
\end{prompt}

\begin{prompt}
\textbf{(c)} Surplus: $\$\answer{6}$
\begin{feedback}
$30-18-6=6$, or $\$2$ each.
\end{feedback}
\end{prompt}
\end{problem}

\begin{problem}
Alex, Bob, and Caroline divide Halloween candy: $4$ Deluxe bars (D), $3$ Walnut Waffles
(W), and $5$ Toffees (T), lined up as follows.
\begin{center}
\begin{tabular}{cccccccccccc}
1 & 2 & 3 & 4 & 5 & 6 & 7 & 8 & 9 & 10 & 11 & 12\\\hline
D & W & D & W & T & W & T & T & D & T & D & T
\end{tabular}
\end{center}
Point values: Alex D $3$, W $2$, T $1$; Bob D $2$, W $3$, T $1$; Caroline D $2$, W $1$,
T $2$. Each player sets markers so that each of their shares is within $1$ point of a third
of their total.

\begin{prompt}
\textbf{(a)} Each player's total points: Alex $\answer{23}$, Bob $\answer{22}$,
Caroline $\answer{21}$
\begin{feedback}
Alex $4\cdot3+3\cdot2+5\cdot1=23$; Bob $4\cdot2+3\cdot3+5=22$; Caroline
$4\cdot2+3+5\cdot2=21$. A third is $7.67$, $7.33$, and $7$.
\end{feedback}
\end{prompt}

\begin{prompt}
\textbf{(b)} One valid set of markers: Alex after pieces $3$ and $8$, Bob after $3$ and
$6$, Caroline after $4$ and $8$. What is each of Bob's shares worth to him?

$1$--$3$: $\answer{7}$ \quad $4$--$6$: $\answer{7}$ \quad $7$--$12$: $\answer{8}$
\begin{feedback}
DWD $=2+3+2=7$; WTW $=3+1+3=7$; TTDTDT $=1+1+2+1+2+1=8$, all within $1$ of $7.33$.
(Alex's shares are worth $8$, $7$, $8$, and Caroline's $6$, $7$, $8$.)
\end{feedback}
\end{prompt}

\begin{prompt}
\textbf{(c)} Alex and Bob tie for the first marker; a coin flip gives Alex the first
share. Complete the allocation.

Alex $1$--$\answer{3}$ \quad Bob $\answer{4}$--$\answer{6}$ \quad Caroline
$\answer{9}$--$12$ \qquad Surplus: pieces $\answer{7}$ and $\answer{8}$
\begin{feedback}
Alex takes $1$--$3$. Of the remaining second markers, Bob's (after $6$) comes before
Caroline's (after $8$), so Bob takes his second share, $4$--$6$. Caroline takes her last
share, $9$--$12$. The two toffees at $7$ and $8$ are left over.
\end{feedback}
\end{prompt}
\end{problem}

\begin{problem}
The country of Alphabitia has four states and uses Hamilton's method to apportion the
$50$ seats of its General Assembly.
\begin{center}
\begin{tabular}{r|rrrrr}
 & A & B & C & D & Total\\\hline
Population & 142 & 27 & 87 & 42 & 298
\end{tabular}
\end{center}

\begin{prompt}
\textbf{(a)} Standard divisor: $\answer{5.96}$
\end{prompt}

\begin{prompt}
\textbf{(b)} Standard quotas (three decimals): A $\answer[tolerance=0.0005]{23.826}$,
B $\answer[tolerance=0.0005]{4.530}$, C $\answer[tolerance=0.0005]{14.597}$,
D $\answer[tolerance=0.0005]{7.047}$
\end{prompt}

\begin{prompt}
\textbf{(c)} Hamilton apportionment: A $\answer{24}$, B $\answer{4}$, C $\answer{15}$,
D $\answer{7}$
\begin{feedback}
Lower quotas total $48$; the $2$ extra seats go to the largest residues, A ($0.826$) and
C ($0.597$).
\end{feedback}
\end{prompt}
\end{problem}

\end{document}
